\section{\bt{Background and Motivation}{背景与动机}}
\label{sec:background}
\subsection{\bt{Application and agent workloads}{应用与智能体工作负载}}
\label{sec:workload-characterization}
\bi{In a parameterized application, query templates and database knowledge reside in code: applications invoke the same queries with different parameters, so millions of queries reduce to thousands of templates~\cite{ma2018querybot}. A fresh data-agent session must acquire that knowledge through schema lookups, data probes, and checks before it can write the query that answers its question. To quantify the difference, we compare two workloads (Table~\ref{tab:workload-characterization}). The application side is Redset~\cite{vanrenen2024redset}, the query metadata of 200 provisioned and 200 serverless Amazon Redshift clusters over three months: 441 million queries, analyzed as complete per-cluster streams, with the dataset's feature-fingerprint template identifier, a proxy that overestimates repetition~\cite{redset2024dataset}. The agent side consists of 100 fresh DeepSeek V4.1 Flash sessions, each answering one of the 25 questions of Section~\ref{sec:evaluation} with only table listing, table description, and SQL tools, three sessions at a time on one database and with no shared state. SQL templates are libpg\_query fingerprints, which ignore constants and formatting.\footnote{\url{https://github.com/pganalyze/libpg_query}} A rule-based classifier, checked on a random sample of calls, labels each call as a schema lookup, data probe, validation, abandoned attempt, final answer, or failure. A lookup's \emph{fact} is its purpose together with the tables and columns it reads, so differently written SQL that probes the same columns counts as the same fact.}
{在参数化应用中，查询模板和数据库知识保存在代码里：应用以不同参数反复调用相同的查询，数百万条查询可归并为数千个模板~\cite{ma2018querybot}。全新的数据智能体会话则必须先通过结构查找、数据探查和检查获得这些知识，才能写出回答问题的查询。为量化这一差别，我们对照两类负载（表~\ref{tab:workload-characterization}）。应用侧采用 Redset~\cite{vanrenen2024redset}，即 Amazon Redshift 中 200 个预置集群和 200 个无服务器集群三个月的查询元数据，共 4.41 亿条查询，按每个集群的完整查询流分析，模板标识采用数据集自带的特征指纹，它是会高估重复程度的近似~\cite{redset2024dataset}。智能体侧是 100 个全新的 DeepSeek V4.1 Flash 会话，每个会话回答第~\ref{sec:evaluation} 节 25 道题中的一道，只有列出表、描述表和执行 SQL 三种工具；三个会话同时访问同一数据库，会话之间不共享状态。SQL 模板使用忽略常量与格式的 libpg\_query 指纹。基于规则的分类器（在随机抽样的调用上人工核对）把每次调用标为结构查找、数据探查、验证、放弃的尝试、最终答案或失败。查找所获得的\emph{事实}由其目的及读取的表和列确定，因此写法不同、但探查相同列的 SQL 算作同一事实。}

\begin{table*}[t]
\caption{\bt{Application and agent workloads. Application: Redset~\cite{vanrenen2024redset}, 441\,M queries from 400 Redshift clusters; pooled rates, with cross-cluster medians where clusters differ widely. Agent: 100 fresh sessions (25 questions, two phrasings, two repetitions, three concurrent) issuing 521 database calls, 176 of them SQL. A \emph{lookup} is a schema lookup, data probe, or validation call.}{应用与智能体负载。应用：Redset~\cite{vanrenen2024redset}，400 个 Redshift 集群的 4.41 亿条查询；给出汇总比例，集群差异大时同时给出跨集群中位数。智能体：100 个全新会话（25 道题、两种题面、各两次，三个并发），共 521 次数据库调用，其中 SQL 176 条。\emph{查找}指结构查找、数据探查或验证调用。}}
\label{tab:workload-characterization}
\footnotesize
\begin{tabularx}{\textwidth}{@{}P{0.15\textwidth}LL@{}}
\toprule
\bt{Dimension}{维度} & \bt{Application (Redset)}{应用（Redset）} & \bt{Data agents (100 sessions)}{数据智能体（100 个会话）}\\
\midrule
\bt{SQL stability}{SQL 稳定性} & \bt{93.4\% of read queries repeat an earlier template (90.1\% of SELECTs); 6.6 templates per 100 queries.}{93.4\% 的读查询重复已出现的模板（SELECT 中 90.1\%）；每 100 条查询 6.6 个模板。} & \bt{24\% of SQL statements repeat an earlier template; 133 templates for 176 statements.}{24\% 的 SQL 重复已出现的模板；176 条 SQL 有 133 个模板。}\\
\bt{Exploration}{探索} & \bt{System-table-only queries: 7.6\% pooled, 2.2\% median.}{只访问系统表的查询：汇总 7.6\%，中位数 2.2\%。} & \bt{68\% of database calls explore, mostly schema lookups; every session starts with one.}{68\% 的数据库调用用于探索，以结构查找为主；每个会话都以结构查找开始。}\\
\bt{Failures}{失败} & \bt{0.35\% of read queries abort (0.54\% of SELECTs).}{0.35\% 的读查询中止（SELECT 中 0.54\%）。} & \bt{1.7\% of SQL fails, but all 32 wrong answers come from SQL that ran without error.}{1.7\% 的 SQL 出错；但 32 个错误答案全部来自执行成功的 SQL。}\\
\bt{Validation}{验证} & \bt{None in parameterized SQL.}{参数化 SQL 中没有。} & \bt{4.8\% of calls; 16\% of sessions check key uniqueness or duplicate keys themselves.}{占调用的 4.8\%；16\% 的会话自行检查键唯一性或重复键。}\\
\bt{Same meaning, different SQL}{同义不同 SQL} & \bt{One template per question type.}{每类题一个模板。} & \bt{87.5\% of questions are answered correctly with $\geq 2$ templates; 47 templates for 15 question types.}{87.5\% 的题目以至少 2 个模板答对；15 类题共 47 个模板。}\\
\bt{Session lifetime}{会话寿命} & \bt{97.8\% of queries come from principals active on $\geq 7$ days; median activity span 80 days.}{97.8\% 的查询来自活跃天数至少 7 天的用户主体；活跃跨度中位数 80 天。} & \bt{Median session lasts 36\,s and makes 5 calls; nothing carries over.}{会话时长中位数 36 秒、5 次调用；不保留任何状态。}\\
\bt{Where DB knowledge lives}{数据库知识的位置} & \bt{In application code.}{在应用代码中。} & \bt{Rediscovered per session: 93\% of lookups repeat a fact an earlier session obtained (77\% of lookup DB time). Most are schema facts; 38\% of SQL lookups repeat a fact, none with repeated SQL text.}{每个会话重新获取：93\% 的查找重复此前会话得到的事实（占查找数据库时间的 77\%）。大部分是表结构事实；SQL 查找中 38\% 重复事实，SQL 文本没有一条重复。}\\
\bt{Concurrent repeated work}{并发重复工作} & \bt{Same template from another user within 60\,s: 0.61\% pooled, 0.01\% median.}{60 秒内其他用户使用同一模板：汇总 0.61\%，中位数 0.01\%。} & \bt{91\% of lookups are repeated by another session within 60\,s.}{91\% 的查找在 60 秒内被其他会话重复。}\\
\bottomrule
\end{tabularx}
\end{table*}


\bi{The two workloads repeat different things. Applications repeat queries: long-lived principals issue the same templates, which prepared statements and result caches serve well (34\% of Redset's read queries are answered from Redshift's result cache); Redshift's own analysis finds that in half of the clusters, 80\% of queries are exact repetitions~\cite{vanrenen2024redset}. Agents repeat knowledge. Their SQL changes from session to session, and correct answers to the same question use different templates. Most repeated lookups are schema facts that a metadata cache could serve; the SQL probes that test business properties repeat less often, and never verbatim, so a cache keyed by query text serves none of them. Two further observations shape the design. Answers are wrong without any failure: 31 of the 32 wrong answers occur when a question names a metric without defining it, exactly where a shared definition helps. And validation is rare: only 16\% of sessions check key uniqueness, the kind of validity condition that \system\ maintains on their behalf. A shared definition removes rediscovery and supplies business meaning, but a reused fact remains correct only while the data it was derived from does not change in the wrong way.}
{两类负载重复的东西不同。应用重复的是查询：长期存在的用户主体发出相同模板，预编译语句和结果缓存能很好地服务它们（Redset 中 34\% 的读查询由 Redshift 的结果缓存直接回答）；Redshift 自身的分析发现，在一半集群中 80\% 的查询是完全相同的重复~\cite{vanrenen2024redset}。智能体重复的是知识。它们的 SQL 在会话之间不断变化，同一道题的正确答案也使用不同模板。重复的查找大部分是表结构事实，元数据缓存可以服务它们；检验业务性质的 SQL 探查重复得少一些，且从不逐字重复，因此按查询文本索引的缓存一条也服务不了。另有两点观察决定了后文的设计。答案在没有任何报错的情况下出错：32 个错误答案中有 31 个出现在题目只给出指标名而不给定义时，这正是共享定义能帮上忙的地方。验证很少见：只有 16\% 的会话检查键唯一性，而这正是 \system\ 代为维护的那类有效性条件。共享定义消除了重复获取并提供业务含义，但被复用的事实只有在其来源数据没有发生错误变化时才保持正确。}

\subsection{\bt{Data agents and metric definitions}{数据智能体与指标定义}}
\bi{A data agent interacts with the database through tools: it lists tables, describes columns, validates joins, and runs SQL, and it submits an answer together with the queries that produced it. A metric definition packages what the agent had to discover. Our running example is store return amount: its definition records the fact table (\code{store\_returns}), the measure (the amount of completed returns), the grain (ticket and item identify one return), the time role (return date), persistent filters, and the convention for empty input; question parameters such as the month are not part of it. Initially, every return has one completed row, so no status filter is needed. A return-rate definition additionally joins sales to returns on the same business key and shares the returns-grain condition with the return-amount definition.}
{数据智能体通过工具与数据库交互：列出表、查看列、验证连接、执行 SQL，并在提交答案时附上产生答案的查询。指标定义把智能体需要摸索出来的内容打包保存。本文以门店退货金额为贯穿例子：其定义记录事实表（\code{store\_returns}）、度量（已完成退货的金额）、粒度（小票号与商品确定一笔退货）、时间角色（退货日）、持久过滤，以及无输入行时的约定；月份等题目参数不属于定义。最初每笔退货只有一条完成记录，因此无需状态过滤。退货率定义还按同一业务键连接销售与退货，并与退货金额定义共享退货粒度条件。}

\bi{Definitions come from two sources. Semantic layers let engineers declare measures, keys, and relationships by hand, and one commercial metric layer documents that a declared join property ``is not validated at runtime'' and that measures return incorrect results if the join fans out~\cite{databricks2026metricviews}. Increasingly, such knowledge is derived from script and query histories~\cite{weng2025datalab} or reused from agent trajectories~\cite{biswal2026agentsm}; learned definitions reach many agents quickly, but their assumptions are recorded, at best, implicitly in the exploration that produced them.}
{定义有两个来源。语义层让工程师手工声明度量、键和关系，而某商用指标层的文档写明，声明的连接性质“不在运行时验证”，一旦连接放大行数，度量就会返回错误结果~\cite{databricks2026metricviews}。越来越多的此类知识则从脚本与查询历史中推导~\cite{weng2025datalab}，或从智能体轨迹中复用~\cite{biswal2026agentsm}；学到的定义能迅速传到许多智能体，但其前提假设至多只隐含在产生它的探索过程中。}

\subsection{\bt{How updates break a correct definition}{更新如何破坏正确的定义}}
\bi{Continue the return-amount example. An ETL update adds an application-state row for each new return while retaining the completed row. The shared SQL, \code{SUM(sr\_return\_amt)}, still runs, but the ticket-and-item key is no longer unique and the aggregate counts some returns twice. Revalidation discovers this broken grain. Adding a completed-status filter is a repair candidate: it restores one row per return while retaining the business keys. The filter must also preserve the stated business meaning and pass regression against learning evidence. The return-rate definition needs the same filtered uniqueness to avoid join fan-out, so one verdict can justify both definitions. And while the update is being loaded, a query that relies on the old verdict must not read the half-changed data as if the grain still held.}
{继续退货金额的例子。一次 ETL 更新在保留完成记录的同时，为新退货追加申请状态行。共享 SQL \code{SUM(sr\_return\_amt)} 仍能执行，但小票号与商品键不再唯一，聚合把部分退货计算了两次。重验证发现粒度被破坏。添加完成状态过滤是一个修复候选：它保留业务键，并恢复每笔退货一条记录。该过滤还必须保留声明的业务含义，并通过学习证据的回归测试。退货率定义也需要同一个过滤后唯一性条件来避免连接放大，因此一个验证结论可以支撑两个定义。而在更新装载期间，依赖旧结论的查询不能把已经变化的数据当作粒度仍然成立来读取。}

\bi{An aggregate is summarizable only if the grouped records are disjoint and complete and the measure is compatible with the aggregation~\cite{lenz1997summarizability}; violations can make analysis tools return wrong results~\cite{mazon2009summarizability}. For a learned definition, disjointness becomes two executable requirements, filtered key uniqueness and bounded join multiplicity, and completeness becomes coverage requirements on business keys and periods. A third requirement, the time role, fixes which date relation assigns a fact to a period. Data-only updates break each of them in practice: status-history rows or duplicate loads violate grain, reused dimension keys make a many-to-one join fan out, and late or re-keyed rows shift a fact to another period or to none. Schema-based invalidation observes none of these changes. When such conditions are checked on data instances rather than guaranteed by the schema, a check holds only for the current data and any change requires a new check~\cite{lenz1997summarizability}. Each requirement can be checked with a SQL probe over a few tables: a grouped count that finds a duplicated key under the filter, or a comparison of row counts before and after a join. Data profiling discovers such unique column combinations while noting that they are valid only for a specific instance~\cite{abedjan2015profiling}, and data-validation systems run such checks declaratively over growing datasets~\cite{schelter2018deequ,breck2019validation}. What is missing for agents is the link between a learned definition and the checks that justify it, so that the checks run when the definition's data changes and their outcome governs whether, and on which data, agents may use the definition.}
{只有当被分组的记录互不相交且完备、度量与聚合方式相容时，汇总才是可汇总的~\cite{lenz1997summarizability}；违反这些条件可能使分析工具给出错误结果~\cite{mazon2009summarizability}。对学到的定义而言，不相交性对应两项可执行要求，即过滤后键唯一性和有界的连接多重性；完备性对应业务键与期间上的覆盖要求。第三项要求是时间角色，它规定由哪种日期关系把事实归入某个期间。纯数据更新在实践中会分别破坏它们：状态流水行或重复装载破坏粒度，被复用的维度键使多对一连接放大行数，迟到或重新赋键的记录把事实移到另一个期间或归不到任何期间。按表结构失效观察不到其中任何一种变化。当这些条件在数据实例上检查、而不由表结构保证时，检查结果只对当前数据成立，任何变化都需要重新检查~\cite{lenz1997summarizability}。每项要求都能用读取少数几张表的 SQL 探查来检查：用分组计数找出过滤后重复的键，或比较连接前后的行数。数据画像能够发现这类唯一列组合，同时指出它们只对某一数据实例成立~\cite{abedjan2015profiling}；数据验证系统能在不断增长的数据集上声明式地执行这类检查~\cite{schelter2018deequ,breck2019validation}。智能体场景中缺少的是把学到的定义与支撑它的检查联系起来：定义所依赖的数据一变，检查就执行，其结果决定智能体能否使用该定义、以及在哪些数据上使用。}

\subsection{\bt{What maintenance must decide}{维护需要决定什么}}
\label{sec:alternatives}
\bi{Table~\ref{tab:ladder} compares maintenance strategies along the three questions of Section~\ref{sec:introduction}. Schema-only maintenance never notices a data-only break. Revoke-on-write never serves a stale definition but discards valid ones after every benign append and depends on LLM relearning. Definition-level lazy revalidation, adapted from lazy view maintenance~\cite{zhou2007lazy}, rechecks all conditions of an affected definition on its next use, and a version-keyed cache of its check queries removes checks repeated across definitions. All of them check before execution and therefore leave a window in which a concurrent write can invalidate the evidence before the query runs. \system\ shares the conditions and repair of definition-level revalidation, reuses verdicts by condition identity and dependency version, and binds each use to the query's own snapshot.}
{表~\ref{tab:ladder} 按第~\ref{sec:introduction} 节的三个问题比较各维护方式。只看结构的维护永远察觉不到纯数据破坏。写入即撤销从不提供过期定义，但每次正常追加后都丢掉有效定义，并依赖大模型重学。借鉴物化视图懒维护~\cite{zhou2007lazy}的定义级懒重验证在受影响定义下次使用时重查其全部条件；按版本缓存其检查查询，可以去掉跨定义重复的检查。它们都在执行前检查，因此留下一个窗口：并发写入可能在查询执行前使证据失效。\system\ 沿用定义级重验证的条件与修复，按条件身份与依赖版本复用结论，并把每次使用绑定到查询自身的快照。}

\begin{table}[t]
\caption{\bt{Maintenance strategies for shared definitions and the studies that compare them: R, library replay; C, controlled cost benchmark; E, end-to-end scenarios (Section~\ref{sec:evaluation}). ``Write during use'': a write that commits between validation and the query that relies on it.}{共享定义的维护方式及比较它们的实验：R，定义库回放；C，受控代价实验；E，端到端情形（第~\ref{sec:evaluation} 节）。“使用期间写入”：在验证与依赖它的查询之间提交的写入。}}
\label{tab:ladder}
\footnotesize
\setlength{\tabcolsep}{3pt}
\begin{tabularx}{\linewidth}{@{}llLLLL@{}}
\toprule
\bt{Strategy}{方式} & \bt{Used in}{实验} & \bt{Benign append}{正常追加} & \bt{Data-only break}{纯数据破坏} & \bt{Write during use}{使用期间写入} & \bt{Work per update}{每次更新的工作}\\
\midrule
\swatch{mSchema}\ \ifbilingual\bhl{Schema-change\\invalidation}{按模式变更失效}\else\bhl{Schema-change\\invalidation}{}\fi & R, E & \bt{keep}{保留} & \bt{stale use}{过期使用} & \bt{unchecked}{不检查} & \bt{none}{无}\\
\swatch{mTable}\ \bhl{Table tests}{表级测试} & R & \bt{keep}{保留} & \bt{invalidate}{失效} & \bt{check, then execute}{先检查后执行} & \bt{declared tests per table}{按表声明的测试}\\
\swatch{mRevoke}\ \ifbilingual\bhl{Invalidate-\\on-write}{写入即失效}\else\bhl{Invalidate-\\on-write}{}\fi & R, E & \bt{relearn}{重学} & \bt{relearn}{重学} & \bt{check, then execute}{先检查后执行} & \bt{LLM relearning}{大模型重学}\\
\swatch{mDef}\ \bhl{Definition-level}{定义级} & R, C, E & \bt{keep}{保留} & \bt{repair}{修复} & \bt{check, then execute}{先检查后执行} & \bt{conditions of affected definitions}{受影响定义的条件}\\
\swatch{mCache}\ \ifbilingual\bhl{Check-result\\cache}{检查结果缓存}\else\bhl{Check-result\\cache}{}\fi & R, C & \bt{keep}{保留} & \bt{repair}{修复} & \bt{check, then execute}{先检查后执行} & \bt{distinct changed checks}{变化的不同检查}\\
\swatch{mCond}\ \system & R, C, E & \bt{keep}{保留} & \bt{bounded repair}{有界修复} & \bt{same snapshot}{同一快照} & \bt{distinct changed conditions}{变化的不同条件}\\
\bottomrule
\end{tabularx}
\end{table}


\bi{Work per update differs because definitions share conditions. After an update of \code{store\_returns} in our controlled workload, definition-level revalidation checks every condition of the eight affected definitions, 19 checks in all, while only three distinct conditions read the updated table. Such sharing is not an artifact of our workload. Treating each aggregate over a fact table in a SELECT block of the 99 TPC-DS query templates as a definition, written independently of this work, yields \TpDefs\ definitions over \TpConds\ distinct conditions, \TpPer\ instances per condition; an update of \code{store\_sales} affects \TpSalesDefs\ definitions with \TpSalesDefChecks\ condition checks but only \TpSalesConds\ distinct conditions. Any mechanism that keys check outcomes by condition and dependency version exploits this sharing, including a generic cache of check queries keyed by SQL text and table versions (Section~\ref{sec:cost}). The harder questions are whether a passed check still holds when the query runs, and what to do when a check fails.}
{每次更新的工作量不同，是因为定义共享条件。在受控负载中更新 \code{store\_returns} 后，定义级重验证要检查 8 个受影响定义的全部条件，共 19 次，而读取更新表的不同条件只有 3 个。这种共享并非本文负载的人为产物。把与本文无关、独立编写的 99 个 TPC-DS 查询模板中每个 SELECT 块对某事实表的聚合视为一个定义，可得到 \TpDefs\ 个定义、\TpConds\ 个不同条件，每个条件平均 \TpPer\ 个实例；更新 \code{store\_sales} 会影响 \TpSalesDefs\ 个定义、\TpSalesDefChecks\ 次条件检查，但不同条件只有 \TpSalesConds\ 个。任何按条件与依赖版本索引检查结果的机制都能利用这种共享，包括按 SQL 文本与表版本索引检查查询的通用缓存（第~\ref{sec:cost} 节）。更难的问题是：通过的检查在查询执行时是否仍然成立，以及检查失败后该怎么办。}

